\section{Results}
\label{sec:results}

Every table body and every number in this section is generated from the per-seed logs by the released analysis script. Unless stated otherwise, values are means over ten seeds on Diabetes 130-US with the MLP backbone; $p$-values are from paired tests across seeds.

\subsection{Main benchmark}
\label{sec:results-fixed}

\begin{table*}[t]
\caption{Main benchmark (Diabetes 130-US, 20 sites, 60 rounds, 10 seeds). Final-round federated AUROC (mean $\pm$ SD), AUPRC, early-stopped AUROC, pooled AUROC, cumulative upload per site, number of transmitting rounds, and one-sided non-inferiority $p$-values against FedPer at $\delta_{\mathrm{NI}}=0.01$ on AUROC and AUPRC.}
\label{tab:main}
\centering\small
\begin{tabular}{lrcccccccc}
\toprule
Method & $n$ & AUROC & AUPRC & ES-AUROC & Pooled & Upload (MB) & Tx rounds & NI$_{\mathrm{AUROC}}$ $p$ & NI$_{\mathrm{AUPRC}}$ $p$\\
\midrule
\gen{main_diabetes_mlp}
\bottomrule
\end{tabular}
\end{table*}

\begin{table}[t]
\caption{Holm-corrected paired comparisons on final AUROC against FedPer (10 seeds): mean difference, BCa 95\% interval, corrected $p$, paired $d_z$.}
\label{tab:holm}
\centering\small
\begin{tabular}{lcccc c}
\toprule
Comparison & $\Delta$ & 95\% CI & $p_{\mathrm{Holm}}$ & $d_z$ & Sig.\\
\midrule
\gen{holm_diabetes_mlp}
\bottomrule
\end{tabular}
\end{table}

Table~\ref{tab:main} reports the main benchmark. Full-rate FedPer reaches \num{diabetesmlpFedperAuroc}~AUROC at \num{diabetesmlpFedperUploadMb}~MB per site. \tierfed{} at $\rho=0.25$ reaches \num{diabetesmlpTierfedAuroc} (paired difference \num{mainTfVsFedperDiff}, $p=\num{mainTfVsFedperP}$) at \num{diabetesmlpTierfedUploadMb}~MB, transmitting in \num{factorialTierfedTxRounds} of 60 rounds; it is \ifflag{diabetesmlpTierfedNonInferiorAuroc}{non-inferior to FedPer on AUROC ($p=\num{diabetesmlpTierfedNiAurocP}$)}{not non-inferior to FedPer on AUROC ($p=\num{diabetesmlpTierfedNiAurocP}$)} and \ifflag{diabetesmlpTierfedNonInferiorAuprc}{on AUPRC ($p=\num{diabetesmlpTierfedNiAuprcP}$)}{not on AUPRC ($p=\num{diabetesmlpTierfedNiAuprcP}$)}. The round-sparse baselines are at least as accurate: AdaComm reaches \num{diabetesmlpAdacommAuroc} at \num{diabetesmlpAdacommUploadMb}~MB (\tierfed{} minus AdaComm: \num{mainTfVsAdacommDiff}, $p=\num{mainTfVsAdacommP}$), and the fixed multi-rate schedule ($K_{\mathrm s}=5$, $K_{\mathrm m}=2$) reaches \num{diabetesmlpFixedAuroc} at \num{diabetesmlpFixedUploadMb}~MB. LAG at its default threshold never skips an upload and is equivalent to full-rate averaging (\num{diabetesmlpLagAuroc}); SCAFFOLD with the Adam local solver used for every method underperforms (\num{diabetesmlpScaffoldAuroc}) and is reported with its native SGD solver in Table~\ref{tab:baselines}. Early-stopped AUROC is within a narrow band for all methods, so no conclusion depends on the final round.

\textbf{Synchronising less does not hurt here.} Table~\ref{tab:periods} sweeps the fixed schedule's periods. A uniform period of four rounds for both shared tiers reaches \num{periodsKsFourKmFourAuroc} at \num{periodsKsFourKmFourUploadMb}~MB, the best accuracy of any configuration in this study; only a long slow-tier period ($K_{\mathrm s}=8$) costs accuracy (\num{periodsKsEightKmOneAuroc} with $K_{\mathrm m}=1$). Table~\ref{tab:factorial} separates the private head from the schedule: FedPer and FedAvg are level (\num{factorialFedperAuroc} and \num{factorialFedavgAuroc}), so the head does not matter on this data, and \tierfed{} with and without the head is also level (\num{factorialTierfedAuroc} and \num{factorialTierfedNoheadAuroc}). The modest gain of every sparse schedule over full-rate averaging is consistent with additional local training acting as personalisation. An earlier version of this work reported that fixed multi-rate schedules lose about 0.05~AUROC to FedPer; that finding was produced by encounters of the same patient appearing in training and test splits, and by a fallback rule that sent the medium tier whenever no tier was due. Both are corrected here and the finding does not reproduce.

\begin{table}[t]
\caption{Fixed-schedule period sweep $(K_{\mathrm s},K_{\mathrm m})$, 10 seeds; rounds with no tier due are local-only.}
\label{tab:periods}
\centering\small
\begin{tabular}{lrcc}
\toprule
$(K_{\mathrm s},K_{\mathrm m})$ & $n$ & AUROC & Upload (MB)\\
\midrule
\gen{periods}
\bottomrule
\end{tabular}
\end{table}

\begin{table}[t]
\caption{Attribution (10 seeds): private head $\times$ schedule, and \tierfed{} with and without the private head.}
\label{tab:factorial}
\centering\small
\begin{tabular}{lrccc}
\toprule
Cell & $n$ & AUROC & AUPRC & Upload (MB)\\
\midrule
\gen{factorial}
\bottomrule
\end{tabular}
\end{table}

\subsection{The budget knob and the controller}
\label{sec:results-budget}

\begin{table}[t]
\caption{Quota sweep for \tierfed{} (10 seeds): accuracy, upload, realised payload fraction, mean intervals chosen by the exact solver, and the share of post-warm-up rounds in which the two shared tiers received different intervals.}
\label{tab:budget}
\centering\small
\resizebox{\columnwidth}{!}{%
\begin{tabular}{lrcccccc c}
\toprule
$\rho$ & $n$ & AUROC & AUPRC & Upload (MB) & Realised & $\bar K_{\mathrm s}$ & $\bar K_{\mathrm m}$ & Differentiated (\%)\\
\midrule
\gen{budget}
\bottomrule
\end{tabular}}
\end{table}

\begin{table}[t]
\caption{Exact solver versus greedy doubling at two quotas (10 seeds).}
\label{tab:solver}
\centering\small
\begin{tabular}{lrcc}
\toprule
$\rho$ / solver & $n$ & AUROC & Upload (MB)\\
\midrule
\gen{solver}
\bottomrule
\end{tabular}
\end{table}

Table~\ref{tab:budget} shows that $\rho$ behaves as an operator knob: upload rises monotonically with the quota, from \num{budgetRhoFifteenUploadMb}~MB at $\rho=0.15$ to \num{budgetRhoEightyUploadMb}~MB at $\rho=0.8$. Accuracy is flat above $\rho=0.25$ and drops only at $\rho=0.15$ (\num{budgetRhoFifteenAuroc}); the best quota is $\rho=0.4$ (\num{budgetRhoFortyAuroc} at \num{budgetRhoFortyUploadMb}~MB). The last column is the honest test of the divergence-driven mechanism: at $\rho=0.25$ the solver gave the two shared tiers different intervals in \num{budgetRhoTwentyFiveDiffPct}\% of rounds, at $\rho=0.4$ in \num{budgetRhoFortyDiffPct}\%, and at most in \num{budgetRhoFifteenDiffPct}\% at $\rho=0.15$. On this data the controller mostly reduces to a uniform period chosen to meet the quota, which explains why a well-chosen fixed period performs as well. Where the solvers differ, exactness helps: at $\rho=0.4$ the exact solver reaches \num{budgetRhoFortyAuroc} at \num{budgetRhoFortyUploadMb}~MB against \num{budgetRhoFortyGreedyAuroc} at \num{budgetRhoFortyGreedyUploadMb}~MB for the greedy doubling rule (Table~\ref{tab:solver}).

\begin{table}[t]
\caption{Baseline sweeps (10 seeds): AdaComm initial period $\tau_0$, LAG threshold $\xi$, and the SGD local solver for FedAvg, FedPer, SCAFFOLD and \tierfed{}.}
\label{tab:baselines}
\centering\small
\begin{tabular}{lrcccc}
\toprule
Configuration & $n$ & AUROC & AUPRC & Upload (MB) & Tx rounds\\
\midrule
\gen{baselines}
\bottomrule
\end{tabular}
\end{table}

Table~\ref{tab:baselines} gives the baselines a fair frontier. AdaComm with $\tau_0=8$ is the most byte-efficient configuration without DP (\num{baselinesAdacommTauEightAuroc} at \num{baselinesAdacommTauEightUploadMb}~MB); LAG skips uploads only at $\xi=10^4$, where its accuracy collapses (\num{baselinesLagXiOneZeroZeroZeroZeroAuroc}). With an SGD local solver SCAFFOLD recovers (\num{baselinesScaffoldSgdAuroc}), FedAvg and FedPer improve (\num{baselinesFedavgSgdAuroc}, \num{baselinesFedperSgdAuroc}), and \tierfed{} is lower (\num{baselinesTierfedSgdAuroc}); the ordering of methods therefore depends on the local solver, which we state as a limitation.

\begin{figure}[t]
\centering\fig{fig_frontier}{0.95\columnwidth}
\caption{Accuracy--communication frontier on Diabetes (10 seeds): every method, quota and compressor on one pair of axes.}
\label{fig:frontier}
\end{figure}

\subsection{Scheduling and compression compose}
\label{sec:results-comp}

\begin{table}[t]
\caption{Scheduling $\times$ compression (10 seeds). Compressors act on the update with per-client error feedback. \tierfed{} rows use $\rho=0.25$.}
\label{tab:comp}
\centering\small
\begin{tabular}{lrccc}
\toprule
Method / compressor & $n$ & AUROC & AUPRC & Upload (MB)\\
\midrule
\gen{compression}
\bottomrule
\end{tabular}
\end{table}

Table~\ref{tab:comp} confirms (\ref{eq:compose}): reductions from timing and from message size multiply. 8-bit quantisation alone takes FedPer from \num{compFedperNoneUploadMb} to \num{compFedperQuantEightUploadMb}~MB at unchanged accuracy (\num{compFedperQuantEightAuroc}); \tierfed{} with 8-bit quantisation uploads \num{compTierfedQuantEightUploadMb}~MB (\num{compTierfedQuantEightAuroc}); and \tierfed{} with sparse ternary compression uploads \num{compTierfedStcUploadMb}~MB at \num{compTierfedStcAuroc}, against \num{compFedperStcUploadMb}~MB (\num{compFedperStcAuroc}) for FedPer with the same compressor. Count-sketch compression fails for every method (\num{compFedavgSketchAuroc}, \num{compFedperSketchAuroc}, \num{compTierfedSketchAuroc}): at the sketch width used, the unsketched estimate of the update is too noisy for this model, and we report it rather than tune it away.

\subsection{Network-level cost}
\label{sec:results-net}

\begin{table}[t]
\caption{Network cost of the main benchmark under the link profiles of Table~\ref{tab:profiles}: total communication time on cell-edge NR, LEO satellite and the 6G edge-slice target, uplink energy on cell-edge NR, and the fraction of rounds fitting a 60-s window on cell-edge NR. Means over 10 seeds, computed from measured bits by (\ref{eq:netcost}).}
\label{tab:network}
\centering\small
\resizebox{\columnwidth}{!}{%
\begin{tabular}{lrrrrrr}
\toprule
Method & Upload (MB) & $t$ NR-edge (s) & $t$ LEO (s) & $t$ 6G (s) & $E$ NR-edge (J) & Window (\%)\\
\midrule
\gen{network_diabetes_mlp}
\bottomrule
\end{tabular}}
\end{table}

Table~\ref{tab:network} translates measured bits into what an operator sees. On a cell-edge NR uplink the total communication time of the job falls from \num{NetCellEdgeFedper}~s for FedPer to \num{NetCellEdgeTierfed}~s for \tierfed{}. The model is small, so every method fits the 60-s window on every link profile; the absolute times matter less than the proportional reductions, which carry over to larger models with the same schedule.

\subsection{A wide model on emulated links}
\label{sec:results-emu}
\begin{table}[t]
\caption{Wide MLP (\num{wideParams} parameters) on Diabetes (3 seeds): accuracy and upload.}
\label{tab:wide}
\centering\small
\begin{tabular}{lrccrr}
\toprule
Method & $n$ & AUROC & AUPRC & Upload (MB) & Params\\
\midrule
\gen{wide}
\bottomrule
\end{tabular}
\end{table}

\begin{table*}[t]
\caption{Trace-driven socket emulation of the wide-MLP runs: total emulated upload time (real TCP over paced loopback links), the analytic model's prediction from (\ref{eq:netcost}), their relative difference, total compute time of the slowest site, and wall-clock time to reach AUROC 0.60 and 0.62 (``never'' if not reached in 60 rounds). Means over seeds; the cell-edge profile was emulated for one seed.}
\label{tab:emu}
\centering\small
\begin{tabular}{llrrrrrrr}
\toprule
Profile & Method & $n$ & Emulated (s) & Analytic (s) & Diff.\ (\%) & Compute (s) & $t_{0.60}$ (s) & $t_{0.62}$ (s)\\
\midrule
\gen{emulation}
\bottomrule
\end{tabular}
\end{table*}

With a model of \num{wideParams} parameters, one full-rate upload per site is several megabytes and the link, not computation, dominates the round on every constrained profile (Table~\ref{tab:emu}). The analytic model agrees with real paced TCP to within \num{emuMaxAbsRelErrPct}\% on every run, so the costs reported in Table~\ref{tab:network} are not an artefact of the formula; the residual is the RTT term, which the analytic model charges in full and the emulator as a one-way delay, plus transport overhead. On the 50-Mbit/s NR profile the emulated upload time is \num{emuNrTypicalFedperUploadS}~s for FedPer and \num{emuNrTypicalTierfedUploadS}~s for \tierfed{}, and the wall-clock time to AUROC 0.60 is \num{emuNrTypicalFedperTtZeropSix}~s against \num{emuNrTypicalTierfedTtZeropSix}~s (AdaComm \num{emuNrTypicalAdacommTtZeropSix}~s, fixed period \num{emuNrTypicalFixedTtZeropSix}~s). Accuracy of the wide model is in Table~\ref{tab:wide}: \tierfed{} \num{wideTierfedAuroc}, FedPer \num{wideFedperAuroc}, fixed period \num{wideFixedAuroc}, AdaComm \num{wideAdacommAuroc}.

\subsection{Attack surface and differential privacy}
\label{sec:results-privacy}

\begin{table}[t]
\caption{Matched-privacy comparison (10 seeds, $C=1$, $\delta=10^{-5}$). Every method is run at the same target $\varepsilon^\star$ with Poisson client sampling at rate $q$; its noise multiplier $\sigma$ is calibrated to its own number of planned releases with the subsampled-Gaussian accountant (\ref{eq:rdp}). Achieved $\varepsilon$, final-round AUROC and AUPRC, and upload. Only data-independent schedules are eligible (Section~\ref{sec:security-dp}).}
\label{tab:dp}
\centering\small
\resizebox{\columnwidth}{!}{%
\begin{tabular}{lrrrrccr}
\toprule
$q$ / $\varepsilon^\star$ / Method & $n$ & $\sigma$ & Releases & $\varepsilon$ & AUROC & AUPRC & Upload (MB)\\
\midrule
\gen{dpmatched}
\bottomrule
\end{tabular}}
\end{table}

\begin{table}[t]
\caption{Sensitivity of \tierfed{}-DP to its prior $\lambda$ in (\ref{eq:jdp}) at $q=1$, $\varepsilon^\star=4$ (10 seeds): chosen interval, calibrated $\sigma$, releases, AUROC, upload.}
\label{tab:dplambda}
\centering\small
\begin{tabular}{lrrrrcr}
\toprule
$\lambda$ & $n$ & $K$ & $\sigma$ & Releases & AUROC & Upload (MB)\\
\midrule
\gen{dplambda}
\bottomrule
\end{tabular}
\end{table}

Table~\ref{tab:dp} compares methods at equal privacy rather than equal noise. A full-rate method makes 60 releases and must use a large $\sigma$ to meet $\varepsilon^\star$; the fixed period makes \num{dpmQHalfEpsFourFixedRel} and \tierfed{}-DP chooses its own release count from (\ref{eq:jdp}) (\num{dpmQHalfEpsFourTierfedDpRel} at $q=0.5$, $\varepsilon^\star=4$, with $\sigma=\num{dpmQHalfEpsFourTierfedDpSigma}$ against \num{dpmQHalfEpsFourFedperSigma} for FedPer). At $q=0.5$ and $\varepsilon^\star=4$, \tierfed{}-DP reaches \num{dpmQHalfEpsFourTierfedDpAuroc} AUROC against \num{dpmQHalfEpsFourFedperAuroc} for FedPer (difference \num{dpmQHalfEpsFourTfdpVsFedperDiff}, $p=\num{dpmQHalfEpsFourTfdpVsFedperP}$), \num{dpmQHalfEpsFourFedavgAuroc} for FedAvg (\num{dpmQHalfEpsFourTfdpVsFedavgDiff}, $p=\num{dpmQHalfEpsFourTfdpVsFedavgP}$) and \num{dpmQHalfEpsFourFixedAuroc} for the fixed period (\num{dpmQHalfEpsFourTfdpVsFixedDiff}, $p=\num{dpmQHalfEpsFourTfdpVsFixedP}$). Across the budget grid, \tierfed{}-DP is \ifflag{dpmQHalfEpsOneTfdpVsFedperBetter}{significantly better than}{not significantly better than} FedPer at $\varepsilon^\star=1$, \ifflag{dpmQHalfEpsTwoTfdpVsFedperBetter}{better}{not significantly better} at $\varepsilon^\star=2$ and \ifflag{dpmQHalfEpsEightTfdpVsFedperBetter}{better}{not significantly better} at $\varepsilon^\star=8$; against the fixed period it is \ifflag{dpmQHalfEpsFourTfdpVsFixedBetter}{better}{\ifflag{dpmQHalfEpsFourTfdpVsFixedWorse}{worse}{level}} at $\varepsilon^\star=4$. Without amplification ($q=1$, $\varepsilon^\star=4$) the corresponding AUROCs are \num{dpmQOneEpsFourTierfedDpAuroc} and \num{dpmQOneEpsFourFedperAuroc}. The planner's choice moves with $\lambda$ as intended (Table~\ref{tab:dplambda}): a larger noise weight selects longer intervals and fewer releases. Because the prior on heterogeneity is flat, \tierfed{}-DP always selects equal intervals for both tiers; its value under DP is that it picks the period, and therefore $\sigma$, from the budget before training, not that it differentiates tiers. The adaptive schedulers are absent from this table by construction: their release patterns depend on client data and are not covered by the accountant. The earlier equal-$\sigma$ comparison that included them is in the Supplement (Section~S3) with that caveat. None of the budgets is claimed to be strong privacy for a 20-site federation.

\begin{table}[t]
\caption{Attack surface (10 seeds): parameters exposed per transmitting round, update-inversion normalised error and feature correlation on the real transmitted slices, membership-inference AUC on the deployed models ($0.5$ = no leakage), and upload.}
\label{tab:attack}
\centering\small
\resizebox{\columnwidth}{!}{%
\begin{tabular}{lrrcccr}
\toprule
Method & $n$ & Params exposed & NMSE $\uparrow$ & Corr.\ $\downarrow$ & MIA AUC $\downarrow$ & Upload (MB)\\
\midrule
\gen{attack}
\bottomrule
\end{tabular}}
\end{table}

Without DP, neither attack separates the methods (Table~\ref{tab:attack}). Update-inversion feature correlation is low for every method (\tierfed{} \num{attackTierfedCorr}, FedPer \num{attackFedperCorr}, AdaComm \num{attackAdacommCorr}), and membership-inference AUC is close to chance (\tierfed{} \num{attackTierfedMia}, FedPer \num{attackFedperMia}). Per transmitted round, \tierfed{} exposes the same slice as FedPer; its reduction in attack surface is only that it transmits in fewer rounds, which this generic attack does not exploit.

\subsection{Robustness, architectures and imaging}
\label{sec:results-arch}

\begin{table}[t]
\caption{Local objective sweep, FedPer versus \tierfed{} ($\rho=0.25$), 10 seeds.}
\label{tab:loss}
\centering\small
\begin{tabular}{lrccc}
\toprule
Loss / method & $n$ & AUROC & AUPRC & Upload (MB)\\
\midrule
\gen{loss}
\bottomrule
\end{tabular}
\end{table}

\begin{table}[t]
\caption{Warm-up sweep for \tierfed{} ($\rho=0.25$), 10 seeds.}
\label{tab:warmup}
\centering\small
\begin{tabular}{lrccc}
\toprule
$T_w$ & $n$ & AUROC & AUPRC & Upload (MB)\\
\midrule
\gen{warmup}
\bottomrule
\end{tabular}
\end{table}

\begin{table}[t]
\caption{Architectures on Diabetes (5 seeds), every method on the same architecture within a row group; parameter count from the instantiated model.}
\label{tab:arch}
\centering\small
\resizebox{\columnwidth}{!}{%
\begin{tabular}{lrcccr}
\toprule
Backbone / method & $n$ & AUROC & AUPRC & Upload (MB) & Params\\
\midrule
\gen{architecture}
\bottomrule
\end{tabular}}
\end{table}

\begin{table}[t]
\caption{PneumoniaMNIST (CNN, 20 sites, Dirichlet $\alpha=0.3$, 10 seeds). Columns as in Table~\ref{tab:main}.}
\label{tab:main-img}
\centering\small
\resizebox{\columnwidth}{!}{%
\begin{tabular}{lrcccccccc}
\toprule
Method & $n$ & AUROC & AUPRC & ES-AUROC & Pooled & Upload (MB) & Tx rounds & NI$_{\mathrm{AUROC}}$ $p$ & NI$_{\mathrm{AUPRC}}$ $p$\\
\midrule
\gen{main_medmnist_cnn}
\bottomrule
\end{tabular}}
\end{table}

\tierfed{} matches or exceeds FedPer under every class-imbalance loss (Table~\ref{tab:loss}; e.g.\ LDAM \num{lossLdamTierfedAuroc} versus \num{lossLdamFedperAuroc}). A longer warm-up buys a small gain at a proportional cost in bytes (Table~\ref{tab:warmup}; $T_w=20$: \num{warmupWarmTwoZeroAuroc} at \num{warmupWarmTwoZeroUploadMb}~MB). On the six-block tabular ResNet, \tierfed{} reaches \num{architectureResnetTierfedAuroc} at \num{architectureResnetTierfedUploadMb}~MB against \num{architectureResnetFedperAuroc} at \num{architectureResnetFedperUploadMb}~MB for FedPer, and on the FT-Transformer \num{architectureTransformerTierfedAuroc} at \num{architectureTransformerTierfedUploadMb}~MB against \num{architectureTransformerFedperAuroc} at \num{architectureTransformerFedperUploadMb}~MB (Table~\ref{tab:arch}). As on the MLP, the fixed multi-rate schedule is at least as accurate as \tierfed{} on every backbone (ResNet \num{architectureResnetFixedAuroc} at \num{architectureResnetFixedUploadMb}~MB; FT-Transformer \num{architectureTransformerFixedAuroc} at \num{architectureTransformerFixedUploadMb}~MB), while \tierfed{} uploads less. The FT-Transformer rows other than the fixed schedule were trained on GPU with the fused attention kernel, which is not bit-deterministic; they are reported over five seeds with that caveat.

On PneumoniaMNIST (Table~\ref{tab:main-img}) the picture differs from the tabular benchmark and we report it as it is. The two methods that share the classifier head are the most accurate, FedAvg (\num{medmnistcnnFedavgAuroc}) and the FedLAMA-style schedule (\num{medmnistcnnFedlamaAuroc}); \tierfed{} reaches \num{medmnistcnnTierfedAuroc} at \num{medmnistcnnTierfedUploadMb}~MB against \num{medmnistcnnFedperAuroc} at \num{medmnistcnnFedperUploadMb}~MB for FedPer, and is \ifflag{medmnistcnnTierfedNonInferiorAuroc}{non-inferior}{not non-inferior} to FedPer at the 0.01 margin ($p=\num{medmnistcnnTierfedNiAurocP}$). Under strong label skew the private head is a liability rather than an asset: both head-private methods trail the head-sharing ones and FedPer has a large between-seed spread. A schedule that keeps the head private should therefore not be used for this kind of task without a shared head.

\subsection{Real institutions}
\label{sec:results-real}
\begin{table}[t]
\caption{Benchmarks with one client per real institution (10 seeds, no DP): Fed-Heart-Disease (4 institutions) and the eICU-CRD demo (\num{realEicuDemoClients} hospitals). AUROC, AUPRC, upload and transmitting rounds.}
\label{tab:real}
\centering\small
\begin{tabular}{lrccrr}
\toprule
Method & $n$ & AUROC & AUPRC & Upload (MB) & Tx rounds\\
\midrule
\multicolumn{6}{l}{\emph{Fed-Heart-Disease}}\\
\gen{realheart}
\midrule
\multicolumn{6}{l}{\emph{eICU-CRD demo}}\\
\gen{realeicudemo}
\bottomrule
\end{tabular}
\end{table}

On the four heart-disease institutions, \tierfed{} reaches \num{realHeartTierfedAuroc} AUROC against \num{realHeartFedperAuroc} for FedPer (difference \num{realHeartTfVsFedperDiff}, $p=\num{realHeartTfVsFedperP}$) and \num{realHeartFedavgAuroc} for FedAvg (\num{realHeartTfVsFedavgDiff}, $p=\num{realHeartTfVsFedavgP}$). On the eICU demo hospitals it reaches \num{realEicuDemoTierfedAuroc} against \num{realEicuDemoFedperAuroc} for FedPer (\num{realEicuDemoTfVsFedperDiff}, $p=\num{realEicuDemoTfVsFedperP}$), the fixed period \num{realEicuDemoFixedAuroc} and AdaComm \num{realEicuDemoAdacommAuroc}. Both federations are small (the heart data has a few hundred patients per institution), so between-seed variance is large and these benchmarks test whether the ordering of methods on the specialty-bucket sites survives real institutional heterogeneity; they are not clinical results. DP was not run on them: with 4 to \num{realEicuDemoClients} clients the client-level Gaussian mechanism needs noise of the order of the whole aggregate.

\subsection{Scalability}

\begin{table}[t]
\caption{Scalability (Diabetes, Dirichlet partition, 30 rounds, 3 seeds): accuracy, upload per site, seconds per round and peak GPU memory, measured, for 10 to 100 sites.}
\label{tab:scaling}
\centering\small
\begin{tabular}{lrcccc}
\toprule
Sites / method & $n$ & AUROC & Upload (MB) & s/round & Peak mem (MB)\\
\midrule
\gen{scaling}
\bottomrule
\end{tabular}
\end{table}

Table~\ref{tab:scaling} reports measured wall-clock and accuracy from 10 to 100 sites with the Dirichlet partition, since the specialty partition cannot produce 100 sites. Upload per site scales linearly with the number of sites for every method and \tierfed{} keeps its proportional saving at every size; seconds per round are dominated by local training and are similar across methods. Accuracy at 100 sites is lower for \tierfed{} than for FedAvg, consistent with each site holding less data. The controller's exact solve enumerates at most 512 points and takes well under a millisecond at every size.
